% Book p. 89 (PDF p. 90)
\section*{Modelling with Algebra}
\subsection*{Application of Algebra}

\subsubsection*{Introduction}

This lesson is continuation of what we learnt in year one functions. We will
explore the properties, graphing techniques and basic theorems of polynomial
functions. Polynomial functions are crucial in mathematics and have applications
in science, engineering and economics. It covers factorising, finding zeros,
graphing, Descartes' Rule of Signs, Fundamental Theorem of Algebra and Linear
and Quadratic Factor Theorems. By the end of the section, you will have a deeper
understanding of polynomial functions and be prepared to solve a variety of
mathematical problems.

\begin{keyideas}
\begin{itemize}
  \item \textbf{Factorising polynomials} involves linear and irreducible quadratic
  factors.
  \item \textbf{Factors} and \textbf{zeros} are variables that make a
  \textbf{polynomial function} equal to \textbf{zero}.
  \item Graphing polynomials with higher degrees reveal complex shapes, key
  features, and roots.
  \item The \textbf{degree} of a polynomial function determines the maximum number
  of solutions a function could have and the number of times a function will
  cross the x-axis when graphed.
  \item The \textbf{degree} of a polynomial function is the total number of
  \textbf{factors}.
\end{itemize}
\end{keyideas}
